% Source: https://artofproblemsolving.com/community/c3h2416069_function_transformations
% Stable identity: aops-topic-2416069-q01
Let $f$ be a function.
\begin{enumerate}
\item[(a)] Show that the graph of $y=f(1-x)$ is obtained from the graph of $y=f(x)$ by reflecting it in the $y$-axis and then translating it $1$ unit to the right, and that this is the same as reflecting it in the line $x=\dfrac12$.
\item[(b)] A student writes $f(1-x)=f\bigl(-(x-1)\bigr)$ and concludes that the graph is obtained by translating $1$ unit to the right and then reflecting in the $y$-axis. Show that this order of steps gives the graph of $y=f(-x-1)$ instead, and explain the mistake.
\item[(c)] For constants $a\ne0$ and $b$, describe how to obtain the graph of $y=f(ax+b)$ from the graph of $y=f(x)$.
\end{enumerate}
